% ── Presentación ─────────────────────────────────────────────────────────────
\section{Presentación}

El presente documento constituye el Plan de Proyecto correspondiente al primer
entregable semestral del curso de Ingeniería de Software, desarrollado por un equipo
de la Escuela Profesional de Ingeniería Informática y de Sistemas de la Universidad
Nacional de San Antonio Abad del Cusco (UNSAAC).

El proyecto tiene como propósito el diseño y desarrollo de una \textbf{Plataforma Web
para la Optimización de la Gestión del TUPA de la UNSAAC}, entendiendo el TUPA
(\emph{Texto Único de Procedimientos Administrativos}) como el instrumento de gestión
institucional que concentra la totalidad de procedimientos administrativos y servicios
prestados por la universidad, junto con sus requisitos, costos, plazos, bases legales
y unidades orgánicas responsables.

Este plan articula la justificación del problema identificado, los objetivos de la
solución propuesta, la metodología de trabajo adoptada, los recursos necesarios, el
cronograma de ejecución, el presupuesto estimado y la descripción del prototipado
dinámico planificado. Su elaboración sigue los principios de la Ingeniería de
Requerimientos y responde a los lineamientos del atributo de graduado AG-C12, con
énfasis particular en la capacidad de adquirir y formular especificaciones de software
de manera sistemática y fundamentada (AG-C12.01).

\begin{supuesto}
\small\textbf{Nota:} Los datos de integrantes y docente deben completarse con la
información real del equipo antes de la entrega formal. Los supuestos asumidos a lo
largo del informe se indican explícitamente en recuadros de este estilo.
\end{supuesto}

